%% kapa: exam=2 round=35 year=2024 period=2 subject=appraisal_theory subject_ko=감정평가이론 seq=2 number=2 points=30
%% source: Q-Net 감정평가사 기출문제 (cst00309) artlSeq=5213532 file=raw/1008/5213532/2024년 제35회 감정평가사 2차시험 문제지.pdf
%% conversion: kordoc 4.24.51 -> tools/gen2.py (auto); review=proofread-sonnet; flags=none
\begin{longquestion}{【문제 2】}
감정평가와 관련된 다음 자료를 참고하여 물음에 답하시오. (30점) \par
\begin{framedtext}
1. 본건은 토지와 건물로 구성된 부동산으로 「집합건물의 소유 및 관리에 관한 법률」 시행 이전에 소유권이전등기가 되어, 현재 ‘건물’은 각 호수별로 등기되어 있고, ‘토지’의 경우도 별도로 등기되어 있음 \par
2. 본건 부동산은 1층(101호, 102호, 103호, 104호, 105호)과 2층(201호, 202호, 203호, 204호, 205호)이 각각 5개호로 구성된 상가로, 현재 건물소유자는 교회 A(101호\textasciitilde{}204호)와 개인 B[205호(교회에 임대됨)]임 \par
3. 상가 전체가 교회로 이용 중이며, 이중 202호, 203호, 204호는 교회의 부속시설로 소예배실, 성경공부방, 교회휴게실로 이용중이고, 용도상 불가분의 관계가 있을 수 있음 \par
4. 202호는 5년 전에, 203호는 3년 전에, 204호는 1년 전에 교회 앞으로 각각 소유권이전등기가 되었고, 건물과 함께 토지 역시 일정 지분이 동시에 교회 앞으로 소유권이전등기됨 \par
5. 건물은 각 호 별로 구조상 독립성과 이용상 독립성이 유지되고 있음 \par
6. 토지는 각 호 별 면적에 비례하여 적정한 지분으로 각 건물소유자들이 공유하고 있음 \par
7. 평가대상 물건은 202호, 203호, 204호이며, 평가목적은 시가참고용임 \par
\end{framedtext}
\medskip\noindent 물음 1) 감정평가사 甲은 평가 대상물건을 개별로 감정평가하기로 결정하였다. 주어진 자료에 근거하여 감정평가사 甲이 개별평가로 결정한 이유를 설명하시오. (10점) \par
\medskip\noindent 물음 2) 반면, 감정평가사 乙은 평가 대상물건을 일괄로 감정평가하기로 결정하였다. 주어진 자료에 근거하여 감정평가사 乙이 일괄평가로 결정한 이유를 설명하시오. (10점) \par
\medskip\noindent 물음 3) 개별평가와 일괄평가의 관점에서 대상물건에 부합하는 평가방법을 설명하시오. (10점) \par
\end{longquestion}
